\documentclass[10pt]{article}
\usepackage[T1]{fontenc}
\usepackage{lmodern,amsmath,amssymb,amsthm,microtype}
\usepackage[margin=0.76in]{geometry}
\usepackage[colorlinks=true,urlcolor=blue,citecolor=blue]{hyperref}
\newtheorem{theorem}{Theorem}
\newtheorem{lemma}{Lemma}
\newtheorem{corollary}{Corollary}
\newcommand{\den}{\operatorname{den}}
\title{Every Denominator Occurs in the Mean of the Divisors\\with Exact Minima at Prime-Power Indices}
\author{Qichao Wang\\{\small Hebei University of Technology}\\[0.3em]
Daoyu Dong\\{\small University of Electronic Science and Technology of China}}
\date{October 26, 2026}
\linespread{0.98}
\hypersetup{pdftitle={Every Denominator Occurs in the Mean of the Divisors},pdfauthor={Qichao Wang and Daoyu Dong}}
\begin{document}
\maketitle
\begin{abstract}
For every positive integer $n$ there is an integer $K$ for which the
reduced denominator of $\sigma(K)/\tau(K)$ is exactly $n$. This proves
the existence conjecture recorded in OEIS A346644. For every odd prime
$q$ and every $s\ge1$, the least such $K$ at denominator $q^s$ is given
by a finite minimum over the partitions of $s$, using the first primes
not congruent to $1$ modulo $q$. The remaining prime-power family has
the closed formula $a(2^s)=2^{2^s-1}$. In particular $a(q)=2^{q-1}$
for every prime $q$.
\end{abstract}

\section{A witness for every denominator}
Let $\sigma(K)$ be the sum of the positive divisors of $K$ and $\tau(K)$
their number. Define
\[
 a(n)=\min\{K\ge1:\den(\sigma(K)/\tau(K))=n\}.
\]
OEIS A346644 \cite{oeis} records this sequence with zero as an escape
value and conjectures that every index has a witness.

\begin{theorem}[Universal denominator construction]\label{thm:exist}
Write $n=2^a m$ with $m$ odd, and put
\[
 u=2^a,\qquad A=2^u-1,\qquad h=mA.
\]
For any odd prime $p\equiv-1\pmod n$, the integer
\begin{equation}\label{eq:K}
 K=2^{u-1}p^{h-1}
\end{equation}
satisfies $\den(\sigma(K)/\tau(K))=n$. Such a prime exists for every $n$.
\end{theorem}
\begin{proof}
Dirichlet's theorem \cite{dirichlet} supplies infinitely many primes
$p\equiv-1\pmod n$, since $\gcd(-1,n)=1$; choose an odd one. For $n=1$
the congruence is vacuous. The factors in \eqref{eq:K} are coprime, so
\[
 \tau(K)=uh=umA=nA,\qquad
 \sigma(K)=A S,\quad S=1+p+\cdots+p^{h-1}.
\]
The integer $h$ is odd. Reducing the alternating sum modulo $n$ gives
$S\equiv1\pmod n$. Hence
\[
 \frac{\sigma(K)}{\tau(K)}=\frac{S}{n},\qquad \gcd(S,n)=1.
\]
This also covers $n=1$, for which $u=A=h=K=1$.
\end{proof}

\section{Exact least values at every prime power}
For a prime $q$, write $v_q$ for its valuation. A factor $p^e$ of $K$
contributes
\[
 c_q(p,e)=v_q(e+1)-v_q(1+p+\cdots+p^e)
\]
to the denominator valuation before taking its positive part.

\begin{lemma}\label{lem:positive}
If $c_q(p,e)>0$, then $p\not\equiv1\pmod q$ and
$c_q(p,e)\le v_q(e+1)$. For any prime $p\not\equiv1\pmod q$ and
integer $b\ge1$,
\begin{equation}\label{eq:replace}
 \den\left(\frac{1+p+\cdots+p^{q^b-1}}{q^b}\right)=q^b.
\end{equation}
\end{lemma}
\begin{proof}
The inequality follows immediately from the nonnegative numerator
valuation. If $q$ is odd and $p\equiv1\pmod q$, lifting the exponent gives
$v_q((p^{e+1}-1)/(p-1))=v_q(e+1)$, so $c_q(p,e)=0$.
For $q=2$ and odd $p$, if $e$ is even both valuations are zero; if $e$
is odd, the elementary $2$-adic lifting formula gives
\[
 v_2((p^{e+1}-1)/(p-1))=v_2(p+1)+v_2(e+1)-1\ge v_2(e+1).
\]
Thus positive contribution again requires $p\not\equiv1\pmod q$.
Finally $p^{q^b}\equiv p\pmod q$ by Fermat's theorem, and $p-1$ is
invertible modulo $q$. The numerator sum in \eqref{eq:replace} is
therefore $1$ modulo $q$, proving the asserted reduced denominator.
\end{proof}

\begin{theorem}[Partition formula]\label{thm:partition}
Fix an odd prime $q$, and list the primes not congruent to $1$ modulo $q$
as $\ell_1(q)<\ell_2(q)<\cdots$. For every $s\ge1$,
\begin{equation}\label{eq:minimum}
 a(q^s)=\min_{\substack{b_1\ge\cdots\ge b_t\ge1\\b_1+\cdots+b_t=s}}
       \prod_{i=1}^t\ell_i(q)^{q^{b_i}-1}.
\end{equation}
The minimum is over the ordinary integer partitions of $s$. For $q=2$,
\begin{equation}\label{eq:two}
 a(2^s)=2^{2^s-1}\qquad(s\ge0).
\end{equation}
Consequently $a(q)=2^{q-1}$ for every prime $q$.
\end{theorem}
\begin{proof}
Every candidate in \eqref{eq:minimum} has divisor count $q^s$ and divisor
sum coprime to $q$, by \eqref{eq:replace}. Its reduced denominator is
exactly $q^s$.

Conversely, let $K=\prod p_i^{e_i}$ have reduced denominator $q^s$.
The sum of all $c_q(p_i,e_i)$ is $s$, so the sum over its positive
terms is at least $s$. By Lemma~\ref{lem:positive}, the corresponding
primes are admissible and their values $a_i=v_q(e_i+1)$ have sum at
least $s$. Choose integers $0\le b_i\le a_i$ with total $s$, discard the
zeros, and replace each exponent by $q^{b_i}-1\le e_i$. The resulting
integer $K'\le K$ has denominator exactly $q^s$ by
\eqref{eq:replace}. Thus it suffices to minimize such products.

For a fixed list of positive $b_i$, use the smallest admissible primes
and assign larger exponents to smaller primes. The latter follows from
$x^Ay^B\le x^By^A$ when $x\le y$ and $A\ge B$. This gives
\eqref{eq:minimum}. There are sufficiently many admissible primes by
Dirichlet's theorem, applied to $-1$ modulo the odd prime $q$.

For $q=2$, the only admissible prime is $2$. The same replacement
argument requires the exponent of $2$ to be at least $2^s-1$.
The integer $2^{2^s-1}$ attains denominator $2^s$, since its divisor
sum is odd; this proves \eqref{eq:two}, including $a(1)=1$.
For odd prime $q$, the only partition of $1$ is $(1)$ and $\ell_1(q)=2$.
\end{proof}

For example $a(9)=\min(2^8,2^2 3^2)=36$ and
$a(25)=\min(2^{24},2^4 3^4)=1296$.
At denominator $81$ the partition $(2,1,1)$ gives
$2^8 3^2 5^2=57600$, smaller than the value $108900$ from $(1,1,1,1)$;
comparison with the other three partitions proves $a(81)=57600$.

The supplied verifier independently computes divisor sums and counts
from integer factorizations, enumerates initial minima, checks the
published values, and evaluates \eqref{eq:minimum}. Large constructive
witnesses are checked modulo their symbolic divisor count, avoiding
expansion of enormous powers.

\begin{thebibliography}{9}
\bibitem{oeis} B. Cloitre, OEIS A346644,
\url{https://oeis.org/A346644}, accessed October 26, 2026.
\bibitem{dirichlet} P. G. L. Dirichlet, \emph{There are infinitely many
prime numbers in all arithmetic progressions with first term and
difference coprime} (1837), translated by R. Stephan,
arXiv:0808.1408v2, \url{https://arxiv.org/abs/0808.1408v2}.
\end{thebibliography}
\end{document}
